\section{Step 106 Build a private validation context}
Fill \texttt{vp\_index\_context}. Construct a fresh index context inside the private stack frame. Copy the caller's directory capacity and scratch address/capacity into its input fields, and initialize its output and retained fields as private storage. Keep the external playback context unchanged.

The same staging address was used before storage and again during reload. An old directory still contains plausible offsets, but it describes an earlier byte sequence. Preparing a new index call establishes that validation will concern the bytes now present at that address.

Use \texttt{s106-1} and stop before \texttt{vp\_index\_context\_ready}.

\paragraph{Observe} The private index context has the supplied directory and scratch controls, with no published playback result.

\Needspace{8\baselineskip}
\paragraph{Check} Why rebuild validation state after reusing the file address?
\begin{enumerate}[label=\Alph*.]
  \item An address alone does not certify the bytes now stored there.
  \item Native pointers change byte order automatically.
  \item Every reload increases the frame count.
\end{enumerate}
